% ==================================================================
% id:       func-021
% topic:    פונקציות
% subtopic: תחום הגדרה ותכונות של פונקציות
% solution: full
% ==================================================================
\textbf{הוכיחו את הטענות הבאות:}
\begin{enumerate}
\item תהי פונקציה $f(x)$ המוגדרת מעל תחום סימטרי $D$. אז קיימות פונקציה זוגית $g(x)$, ופונקציה אי זוגית $h(x)$, כך שמתקיים $f(x)=g(x)+h(x)$.
\hint{חפשו את $g$ ואת $h$ כצירופים של $f(x)$ ושל $f(-x)$.}
\begin{solbox}
בשאלה \qref{func-007} ראיתם כי
$g(x)=\frac{f(x)+f(-x)}{2}$
זוגית\\

ו-
$h(x)=\frac{f(x)-f(-x)}{2}$
אי זוגית.
מתקיים
$f(x)=g(x)+h(x)$
\end{solbox}
\item תהי פונקציה $f(x)$ שהיא גם זוגית וגם אי זוגית אז $f(x)=0$.
\hint{כתבו את $f(-x)$ בשתי דרכים: פעם לפי הזוגיות ופעם לפי האי-זוגיות.}
\begin{solbox}
$f(-x)=f(x)$
מכך שזוגית
$f(-x)=-f(x)$
מכך שאי זוגית ולכן
\\
$-f(x)=f(-x)=f(x)$
לכן
$f(x)=0$
.
\end{solbox}
\item תהי פונקציה $f(x)$ זוגית מצאו עבורה פירוק כמו בסעיף א'.
\begin{solbox}
$f(x)=g(x)+0$
הרי 0  אי זוגית
\end{solbox}
\item תהי פונקציה $f(x)$ אי זוגית מצאו עבורה פירוק כמו בסעיף א'.
\begin{solbox}
$f(x)=h(x)+0$
הרי 0   זוגית
\end{solbox}
\item הראו כי הפירוק בסעיף א' הוא יחיד.
כלומר אם קיימות פונקציה זוגית $g(x)$
ופונקציה אי זוגית $h(x)$
כך שמתקיים $f(x)=g(x)+h(x)$,
וכן קיימות
פונקציה זוגית $g'(x)$
ופונקציה אי זוגית $h'(x)$ \\
כך שמתקיים $f(x)=g'(x)+h'(x)$.
אז $f(x)=f'(x),g(x)=g'(x)$
.
\hint{העבירו אגפים כך שבצד אחד תתקבל פונקציה זוגית ובצד השני פונקציה אי-זוגית, והשתמשו בסעיף ב'.}
\begin{solbox}
$f(x)=g(x)+h(x)=g'(x)+h'(x)$
נתבונן בהפרש
$g(x)-g'(x)=h'(x)-h(x)$
באגף שמאל יש הפרש של פונקציות זוגיות ולכן זוגי ובאגף ימין יש הפרש של פונקציות אי זוגיות ולכן אי זוגי.
מכאן לפי סעיף ב' שני האגפים שווים ל-0
ולכן הפירוק יחיד.
\end{solbox}
\end{enumerate}
